\documentclass[11pt,a4paper]{article}
\usepackage[T1]{fontenc}
\usepackage{lmodern}
\usepackage[margin=22mm]{geometry}
\usepackage{amsmath,amssymb,amsthm,booktabs}
\usepackage[hidelinks]{hyperref}
\hypersetup{pdftitle={Total dominator colorings of hypercubes and total co-independent domination in prisms},
            pdfauthor={},
            pdfsubject={Research preprint, prepared for contributor review}}
\newtheorem{theorem}{Theorem}
\newtheorem{lemma}[theorem]{Lemma}
\newtheorem{corollary}[theorem]{Corollary}
\newcommand{\coi}{\gamma_{t,\mathrm{coi}}}
\newcommand{\boxprod}{\mathbin{\square}}
\newcommand{\td}{\chi_{\mathrm{td}}}
\newcommand{\F}{\mathbb F}
\newcommand{\wt}{\operatorname{wt}}
\title{Total dominator colorings of hypercubes and total co-independent domination in prisms}
% Enter the credited human contributor(s) before a public deposit.
\author{}
\date{Working draft, 27 November 2026}

\begin{document}
\maketitle
\begin{abstract}
We prove that the total dominator chromatic number of the binary
hypercube $Q_d$ is $\gamma_t(Q_d)+2$ for every $d\ge4$, where $\gamma_t$
denotes the total domination number. This extends the dimension families
treated by Kazemi and by Gebendorfer. The proof combines Gebendorfer's
color-class count with bounds on radius-one covering codes and on
total dominating vertex covers. For a prism over a graph $H$ without
isolated vertices, we show that the minimum size of a total dominating
vertex cover is $|V(H)|$ plus the least size of a dominating set $S$
such that $H-S$ is bipartite. As a consequence, a question of Alfred Raju M.,
Palagiri and Yero about equality in their prism upper bound has an
affirmative answer for bases without isolated vertices. Allowing isolated
vertices gives a counterexample on five vertices.
\end{abstract}

\section{The coloring problem}

The vertices of $Q_d$ are the binary words of length $d$, with two words
adjacent when they differ in one coordinate. A set $S$ is \emph{total
dominating} if every vertex, including every vertex of $S$, has a neighbor in
$S$. Write $g_d=\gamma_t(Q_d)$ for the minimum size of such a set. A
\emph{total dominator coloring} is a proper coloring in which every vertex is
adjacent to all vertices of at least one color class. Write $t_d=\td(Q_d)$
for the least number of colors.

Kazemi proved $g_d\le t_d\le g_d+2$ and determined $t_d$ for $d\le4$
\cite[Propositions 4.3 and 4.9]{kazemi14}. In 2026 he proved that the
upper bound is attained in dimensions $2^n$ for $n\ge3$ and conjectured
that it is attained in dimensions $2^n+1$ for $n\ge2$
\cite[Theorem 5.1 and Conjecture 6.1]{kazemi26}.
Gebendorfer proved that conjecture in October 2026
\cite[Theorem 4.3]{geben}. The known total domination number
\cite[Corollary 13]{prisms}, also recorded in \cite[Theorem 1.3]{kazemi26}, gives
$t_{2^n+1}=2^{2^n-n+1}+2$.

\begin{theorem}\label{thm:main}
For every integer $d\ge4$, one has $\td(Q_d)=\gamma_t(Q_d)+2$.
\end{theorem}

For a general dimension, a coloring with at most $g_d+1$ colors has at
most one color class with no common neighbor, by Gebendorfer's argument.
If there is such a class, its complement is both a total dominating set
and a vertex cover. We first determine the smallest size of such a
complement on a prism. This supplies the lower bound in the coloring
problem, including the exceptional estimate needed for $d=6$.
The earlier prediction $t_d=2^{d-2}+2$ for all $d\ge5$
\cite[Conjecture 4.10]{kazemi14} already fails for $d=6$.
Indeed, $g_6=14$ by the prism identity \cite{prisms} and the known
domination number of $Q_5$ \cite[Table 1.1]{kazemi26}, so even the
previous upper bound is 16. Theorem~\ref{thm:main} gives equality.

\section{Total co-independent domination in prisms}

A \emph{total co-independent dominating set} is a total dominating set
whose complement is a nonempty independent set. Write $\coi(G)$ for its
minimum size. The ordinary domination number $\gamma(H)$ is the minimum
size of a set containing a neighbor of each vertex outside it. The prism $H\boxprod K_2$ has two layers indexed by $0,1$ and a
vertical edge between the two copies of each vertex of $H$.

For graphs $H$ without isolated vertices, Alfred Raju M., Palagiri and
Yero proved the following bounds, stated in terms of the covering total
Italian domination number. Their Theorem~2.3 identifies that parameter
with $\coi$ when the minimum degree is at least two, as it is in the prism
\cite[Theorems 2.3 and 5.1]{italian}. If $H$ has order $m$ and
independence number $\alpha(H)$, their bounds are
\begin{equation}\label{eq:oldprism}
 m+\gamma(H)\le\coi(H\boxprod K_2)\le 2m-\alpha(H).
\end{equation}
For paths, Dutton and Klostermeyer determined the total vertex cover
number of the ladder and observed that neighboring fibers with one
selected vertex must use opposite layers \cite[Proposition 30]{dutton}.
The following formula keeps track of this choice of layer in the
lower-bound proof of \eqref{eq:oldprism}.

\begin{theorem}\label{thm:prism}
Let $H$ be a nonempty finite graph without isolated vertices and put
$m=|V(H)|$. Then
\[
  \coi(H\boxprod K_2)
  =m+\min\{|S|: S\text{ dominates }H\text{ and }H-S\text{ is bipartite}\}.
\]
In particular, for bipartite $H$ this is $m+\gamma(H)$.
\end{theorem}

\begin{proof}
Let $T$ be a total co-independent dominating set of the prism. It covers
each vertical edge. Let $S$ consist of those $u$ for which both copies
belong to $T$. Thus $|T|=m+|S|$. We recall the double-fiber observation
from the proof of \cite[Theorem 5.1]{italian}. For $u\notin S$, let
$(u,i)$ be its unique copy in $T$. Its neighbor in $T$ must be a
horizontal neighbor $(v,i)$. The edge from $(u,1-i)$ to $(v,1-i)$
forces $(v,1-i)\in T$ as well, hence $v\in S$. This proves that $S$
dominates $H$. Moreover, the layer of the omitted copy of each vertex
of $H-S$ gives a proper two-coloring, since omitted vertices are
independent.

Conversely, take a minimizing $S$ and a proper two-coloring of $H-S$.
Keep both copies of each vertex in $S$, and omit just the copy in
the assigned layer for each vertex of $H-S$. The omitted vertices
are independent. Every vertex over $S$ has its vertical mate in the set.
Every omitted vertex also has its mate in the set. Every remaining vertex
over $u\notin S$ has a horizontal neighbor over $S$ by domination.
The complement is nonempty since deleting all but any one vertex of $H$
gives an admissible choice of size $m-1$. This proves the reverse bound.
\end{proof}

For bipartite $H$, equality in the lower bound of \eqref{eq:oldprism}
also follows by combining earlier results. Applied to
$X=H\boxprod K_2$, the bipartite bound of Mojdeh et al.\ gives
$\coi(X)\le (|V(X)|+\gamma_t(X))/2$ since $\delta(X)\ge2$
\cite[Theorem 4.10 in the preprint]{mojdeh}. The bipartite prism
identity of Azarija, Henning and Klav\v{z}ar gives
$\gamma_t(X)=2\gamma(H)$ \cite{prisms}, and \eqref{eq:oldprism}
provides the matching lower bound.

Thus for $d\ge2$,
\begin{equation}\label{eq:coicube}
  \coi(Q_d)=2^{d-1}+\gamma(Q_{d-1})=(2^d+g_d)/2.
\end{equation}
The second equality also has a direct verification. Deleting the last
coordinate takes each
parity side of $Q_d$ bijectively onto $Q_{d-1}$ and maps each open
neighborhood onto a closed ball of radius one. Thus a total dominating
set projects on each side to a dominating set of $Q_{d-1}$, and using a
minimum dominating set on each side gives equality.

Formula~\eqref{eq:coicube} gives a proof of the equality that Sonkamble,
Patekar and Waphare state both in Theorem~3.1 and as a conjecture on the
same page \cite[p.~460]{sonkamble}. For $d=2^r$ with $r\ge1$, the classical binary
Hamming code covers $Q_{d-1}$ with $2^{d-1}/d$ words. Its parity-check
matrix consists of the nonzero $r$-bit columns, so every nonzero syndrome
can be corrected by one bit flip. No smaller code can cover, since each
closed ball has $d$ words. Hence $\coi(Q_d)=2^{d-1}+2^{d-1}/d$.

\section{Counting vertices and common neighbors}

Use the notation $N(v)$ for the open neighborhood of $v$. For a nonempty
color class $C$, let $B(C)=\bigcap_{v\in C}N(v)$, and call $C$ active
if $B(C)$ is nonempty. The open neighborhoods of two distinct cube
vertices intersect only when the vertices are at distance two, in which
case their intersection has size two. This observation underlies Kazemi's
classification of the possible common neighborhoods
\cite[proof of Proposition 4.2]{kazemi14}, \cite[Fact 5.1]{kazemi26}.
Gebendorfer expressed it in the following form for every active class and
$d\ge3$ \cite[Lemma 4.1]{geben}.
\begin{equation}\label{eq:local}
                   |C|+|B(C)|\le d+1.
\end{equation}
Indeed, if either set is a singleton, the other has at most $d$ elements.
If both have at least two elements, all edges between the two sets are
present, so each set has at most two elements.

\begin{lemma}\label{lem:count}
Let $d\ge4$, $N=2^d$, and $g=g_d$. A total dominator coloring of $Q_d$
with at most $g+1$ colors implies at least one of the two inequalities
\begin{align}
                  (d+1)(g+1)&\ge2N,\label{eq:all}\\
                       (2d+1)g&\ge3N.\label{eq:one}
\end{align}
In particular, $(d+1)g\le 3N/2$ implies $t_d=g+2$.
\end{lemma}

\begin{proof}
We follow the two cases in Gebendorfer's proof of Theorem~4.3
\cite{geben}. The sets $B(C)$ cover the vertex set. Selecting one
representative from each active class gives a total dominating set,
so there are at least $g$ active classes \cite[Lemma 4.2]{geben}. If every class is active, summing \eqref{eq:local}
counts all $N$ vertices as members of classes and at least $N$ vertices as
common neighbors. This gives \eqref{eq:all}.

Otherwise there are exactly $g$ active classes and one inactive class $D$.
Put $W=V(Q_d)\setminus D$. Every vertex has a neighbor among the
representatives of active classes, all of which lie in $W$. As $D$ is
independent, $W$ is a total co-independent dominating set. A general
degree bound of Cabrera-Mart\'inez et al.\ \cite[Theorem 8, $k=1$]{outer}
gives $|W|\ge dN/(2d-1)$. Here \eqref{eq:coicube} gives
$|W|\ge (N+g)/2$. Summing \eqref{eq:local} over the active classes
yields $(d+1)g\ge N+|W|\ge(3N+g)/2$, proving \eqref{eq:one}.

If $(d+1)g\le3N/2$, then $(d+1)(g+1)<2N$ since $d+1<N/2$,
and $(2d+1)g\le3N-g<3N$. For the matching upper bound,
give every vertex of a minimum total dominating set its own color and
color the remaining vertices by parity using two additional colors.
Every vertex sees a singleton class.
\end{proof}

\section{The dimension estimate}

Write $K(r,1)=\gamma(Q_r)$ for the smallest number of binary
words whose closed radius-one balls cover $Q_r$. The prism identity
above reads $g_{r+1}=2K(r,1)$.

Set $p_d=g_d/2^d$ and $F(d)=(d+1)p_d$. A total dominating set remains
total dominating after a free binary coordinate is appended to its
vertices. For doubling the dimension we use the classical construction
of Mollard and of Katsman and Litsyn, as recorded by Cohen, Lobstein
and Sloane \cite[Theorem 13]{cls}. In covering-code notation it gives
$K(2r+1,1)\le2^rK(r,1)$. Here is the corresponding construction for
total domination. If $S$ totally dominates $Q_a$, put
\[
  T=\{(x,y)\in \F_2^a\times\F_2^a:
           x+y\in S,\ \wt(x)\equiv0\pmod 2\}.
\]
For each $s\in S$ and each even-weight $x$ there is exactly one choice of
$y$, so $|T|=2^{a-1}|S|$. Given $(x,y)$, choose a unit vector $e_i$
such that $x+y+e_i\in S$. If $x$ has odd weight, flip $x_i$. If $x$
has even weight, flip $y_i$. The resulting neighbor belongs to $T$ in
either case. Consequently
\begin{equation}\label{eq:double}
  p_{2a},\ p_{2a+1}\le \tfrac12p_a,
  \qquad F(2a),\ F(2a+1)\le F(a).
\end{equation}

\begin{proof}[Proof of Theorem~\ref{thm:main}]
The two words $000,111$ cover $Q_3$, giving $p_4\le1/4$ and
$p_5\le1/4$ by free extension. The bounds $K(5,1)\le7$ and
$K(6,1)\le12$ give $p_6\le7/32$ and $p_7\le12/64$.
The short codes are recorded below. Thus $F(d)\le3/2$ for $d=4,5,7$.

For $d=6$, $g=2K(5,1)\le14$.
Then $7(g+1)\le105<128$ and $13g\le182<192$. Neither alternative in
Lemma~\ref{lem:count} can hold, giving $t_6=g+2$.

In dimension $12$, \eqref{eq:double} gives $F(12)\le13\cdot7/64<3/2$.
In dimension $13$ we use Wille's 1996 bound $K(12,1)\le380$ \cite{wille}.
It gives $F(13)\le14\cdot380/4096=665/512<3/2$.
For $d\ge8$ other than $12$ and $13$, the index
$a=\lfloor d/2\rfloor$ is at least four and differs from six.
Strong induction using \eqref{eq:double} gives $F(d)\le3/2$ for every
$d\ge4$ except six. Lemma~\ref{lem:count} finishes the proof.
\end{proof}

\section{Equality in the prism upper bound}

After proving \eqref{eq:oldprism}, Alfred Raju M., Palagiri and Yero
asked whether attaining its upper bound forces $\gamma(H)=m-\alpha(H)$
\cite[Section 5.1, p.~341]{italian}. For bases without isolated vertices,
the answer is affirmative. The characterization of graphs whose domination and
vertex cover numbers coincide is classical. Lingas et al.\ describe the
role of leaf neighbors of a minimum vertex cover under that equality
\cite[Theorem 1]{lingas}. We use this property under the prism equality below.

\begin{corollary}\label{cor:upper}
Under the hypotheses of Theorem~\ref{thm:prism},
$\coi(H\boxprod K_2)=2m-\alpha(H)$ if and only if
$\gamma(H)=m-\alpha(H)$.
\end{corollary}

\begin{proof}
If $\gamma(H)=m-\alpha(H)$, the two bounds in \eqref{eq:oldprism}
coincide. Conversely, assume equality in the upper bound. Choose a
maximum independent set $I$ and let $C=V(H)\setminus I$. By
Theorem~\ref{thm:prism}, the smallest dominating set whose deletion
leaves a bipartite graph has size $|C|$.

Suppose $x\in C$ has a neighbor in $C$. If $x$ had no degree-one
neighbor in $I$, then $C\setminus\{x\}$ would still dominate $H$.
Indeed, $x$ has a neighbor in this set. Each vertex of $I$ has a neighbor
in $C$, and only a leaf joined to $x$ could lose its last such neighbor.
The complement
$I\cup\{x\}$ is bipartite, contradicting minimality of $|C|$.
Thus each vertex incident with an edge of $H[C]$ has a leaf neighbor in $I$.

A minimum dominating set $S$ of $H$ can be chosen to contain all these
vertices. Whenever such a vertex is missing, its leaf neighbor must lie in
$S$, and replacing that leaf by its neighbor preserves domination.
The replacements discard only vertices of $I$. Now both $I\setminus S$
and $C\setminus S$ are independent, so $H-S$ is bipartite. The minimum
in Theorem~\ref{thm:prism} is therefore $\gamma(H)$, proving
$\gamma(H)=|C|=m-\alpha(H)$.
\end{proof}

If the question for the covering total Italian domination number is read
as allowing isolated vertices in the base, take $H=K_4\sqcup K_1$.
Its prism has covering total Italian domination number $6+2=8$.
For clarity, an admissible labeling assigns $0,1,2$ to the vertices.
The zeros must be independent and have neighbor-label sum at least two,
and the positive vertices must totally dominate. In $K_4\boxprod K_2$
each layer has at most one zero, so weight at least six is required.
Putting a zero in each layer in different fibers and assigning one to
the other six vertices attains it. Both vertices of the $K_2$ component
require a positive neighbor, giving weight two. Consequently the optimum weight is
$2m-\alpha(H)=8$, although
$\gamma(H)=2<3=m-\alpha(H)$. In fact the minimization in
Theorem~\ref{thm:prism} still computes the covering total Italian
domination number when the base has isolates. Each isolated base vertex
adds one to $m$, must lie in every dominating set, and contributes a
$K_2$ of optimum weight two to the prism. If the base consists entirely of
isolates, the same count gives weight $2m$.

\section*{Short covering certificates}

The covering numbers at lengths five and six are classical
\cite{cls,kazemi26}. The following sets give short certificates for the
upper bounds used in the proof.

For $r=5,6$, write a word as a prefix $u$ and a three-bit suffix. Encode
both as integers, with leading zeroes understood. The entry $S_u$ lists
the suffixes of the codewords with prefix $u$. The entry $U_u$ lists the
suffixes outside the closed radius-one balls of $S_u$ in $Q_3$. A string of
digits in the table lists a set of individual digits.
\begin{center}
\begin{tabular}{cc@{\qquad}rrrrrrrr}
\toprule
$r$ && $u=0$ & 1 & 2 & 3 & 4 & 5 & 6 & 7\\
\midrule
5 & $S_u$ & $0$ & $567$  & $34$ & $0$ \\
  & $U_u$ & $3567$ & $0$ & $\varnothing$ & $3567$ \\
\midrule
6 & $S_u$ & $03$ & $5$ & $16$ & $2$ & $5$ & $16$ & $2$ & $47$ \\
  & $U_u$ & $56$ & $0236$ & $\varnothing$ & $1457$ &
                  $0236$ & $\varnothing$ & $1457$ & $12$ \\
\bottomrule
\end{tabular}
\end{center}
In each column every digit in $U_u$ occurs in an $S_v$ whose prefix
$v$ differs from $u$ in one bit. This verifies coverage. Counting the
digits in the two $S_u$ rows gives seven and twelve codewords.

\paragraph{AI disclosure.}
OpenAI language models derived the extensions presented here, searched the
literature and drafted the manuscript. Separate model instances checked
the proofs. Computations verified the short covering certificates and
the prism results on all 814 labelled graphs of orders two through five
without isolated vertices, including exhaustive checks of the ternary
labels on their prisms. No independent human review has been recorded
for this version.

\begin{thebibliography}{99}
\raggedright

\bibitem{prisms}
J.~Azarija, M.~A.~Henning and S.~Klav\v{z}ar,
(Total) domination in prisms,
\emph{Electron. J. Combin.} \textbf{24} (2017), P1.19.
\href{https://doi.org/10.37236/6288}{doi:10.37236/6288}.

\bibitem{outer}
A.~Cabrera-Mart\'inez, J.~C.~Hern\'andez-G\'omez, E.~Parra-Inza and
J.~M.~Sigarreta Almira, On the total outer $k$-independent domination
number of graphs, \emph{Mathematics} \textbf{8} (2020), 194.
\href{https://doi.org/10.3390/math8020194}{doi:10.3390/math8020194}.

\bibitem{cls}
G.~D.~Cohen, A.~C.~Lobstein and N.~J.~A.~Sloane,
Further results on the covering radius of codes,
\emph{IEEE Trans. Inform. Theory} \textbf{32} (1986), 680--694.
\href{https://doi.org/10.1109/TIT.1986.1057227}{doi:10.1109/TIT.1986.1057227}.

\bibitem{dutton}
R.~Dutton and W.~F.~Klostermeyer, Edge dominating sets and vertex covers,
\emph{Discuss. Math. Graph Theory} \textbf{33} (2013), 437--456.
\href{https://doi.org/10.7151/dmgt.1681}{doi:10.7151/dmgt.1681}.

\bibitem{geben}
J.~J.~Gebendorfer, Total dominator colorings of hypercubes of dimension
$2^n+1$, preprint, 14 October 2026.
\href{https://doi.org/10.13140/RG.2.2.31321.28006}{doi:10.13140/RG.2.2.31321.28006}.
\href{https://www.researchgate.net/publication/415823778_Total_dominator_colorings_of_hypercubes_of_dimension_2_n_1}{Preprint}.

\bibitem{kazemi14}
A.~P.~Kazemi, Total dominator coloring in product graphs,
\emph{Util. Math.} \textbf{94} (2014), 329--345.
\href{https://combinatorialpress.com/um/vol94/}{Journal volume}.

\bibitem{kazemi26}
A.~P.~Kazemi, Total dominator coloring of the hypercube $Q_{2^n}$,
\emph{Contrib. Math.} \textbf{13} (2026), 1--10.
\href{https://doi.org/10.47443/cm.2025.102}{doi:10.47443/cm.2025.102}.

\bibitem{lingas}
A.~Lingas, M.~Miotk, J.~Topp and P.~\.{Z}yli\'nski,
Graphs with equal domination and covering numbers,
\emph{J. Comb. Optim.} \textbf{39} (2020), 55--71.
\href{https://doi.org/10.1007/s10878-019-00454-6}{doi:10.1007/s10878-019-00454-6}.

\bibitem{mojdeh}
D.~A.~Mojdeh, I.~Peterin, B.~Samadi and I.~G.~Yero,
On three outer-independent domination related parameters in graphs,
\emph{Discrete Appl. Math.} \textbf{294} (2021), 115--124.
\href{https://doi.org/10.1016/j.dam.2021.01.027}{doi:10.1016/j.dam.2021.01.027}.
\href{https://arxiv.org/abs/1812.10946}{Preprint}.

\bibitem{italian}
Alfred Raju M., V.~S.~R.~Palagiri and I.~G.~Yero,
On the total version of the covering Italian domination problem,
\emph{Discrete Appl. Math.} \textbf{358} (2024), 333--343.
\href{https://doi.org/10.1016/j.dam.2024.07.017}{doi:10.1016/j.dam.2024.07.017}.

\bibitem{sonkamble}
R.~A.~Sonkamble, S.~C.~Patekar and B.~N.~Waphare,
Total co-independent domination in Cartesian product of graphs with
hypercubes and its variants, \emph{Adv. Appl. Discrete Math.}
\textbf{43} (2026), 451--467.
\href{https://doi.org/10.17654/0974165826029}{doi:10.17654/0974165826029}.

\bibitem{wille}
L.~T.~Wille, New binary covering codes obtained by simulated annealing,
\emph{IEEE Trans. Inform. Theory} \textbf{42} (1996), 300--302.
\href{https://doi.org/10.1109/18.481808}{doi:10.1109/18.481808}.

\end{thebibliography}
\end{document}
